\section{On the decomposition of $\Xfv$ into $K$-orbits}

\label{section-3}

The purpose of this section is to prove the results stated in Theorems \ref{Tp1} and \ref{C1}, regarding the decomposition of $\Xfv=\Grass(V,r)\times\Flags(V^+)\times\Flags(V^-)$ into orbits of $K=\GL(V^+)\times\GL(V^-)$.

By $B_k^+\subset\GL_k(\mC)$ we denote the subgroup of invertible upper triangular matrices in $\GL_k(\mC)$.
Then\[B_K=\left\{\begin{pmatrix}
a & 0\\ 0 & d
\end{pmatrix}:a\in B_p^+,\ d\in B_q^+\right\}=B_p^+\times B_q^+\subset\GL(V^+)\times\GL(V^-)\]
is a Borel subgroup of $K$. Recall that $\flag_0^+$ and $\flag_0^-$ denote the standard flags of $V^+$ and $V^-$. Thus, $B_K$ is the stabilizer of the pair $(\flag_0^+,\flag_0^-)$ for the action of $K$ on the product of flag varieties $\Flags(V^+)\times\Flags(V^-)$.
In Section \ref{section-3.1}, we observe that there is a one-to-one correspondence between the orbit sets $\Xfv/K$ and $\Grass(V,r)/B_K$, which preserves the inclusion relations between closures of orbits. This fact is a useful ingredient in the rest of the section.

In Section \ref{section-3.2}, we show the parametrization of orbits and the dimension formula stated in Theorem \ref{Tp1}\ref{Tp11}--\ref{Tp13}. In Section \ref{section-3.3}, we describe the closure relations of orbits
by proving Theorems \ref{Tp1}\ref{Tp14} and \ref{C1}.
In Section \ref{section-3.4}, we make further remarks and mention relations with the existing literature.


\subsection{A preliminary lemma}

\label{section-3.1}

We will use the following lemma.


\begin{lemma}
\label{L3.1}
\begin{enumerate}
\item The mapping $\Grass(V,r)\to\Xfv$, $W\mapsto (W,\flag_0^+,\flag_0^-)$ induces a one-to-one correspondence between the orbit sets
\[\Xi:\Grass(V,r)/B_K\to \Xfv/K,\ \Gorbit=B_K\cdot W\mapsto \Xi(\Gorbit)=K\cdot (W,\flag_0^+,\flag_0^-).\]
\item If $\Xorbit=\Xi(\Gorbit)\subset\Xfv$ is the $K$-orbit corresponding to $\Gorbit\subset\Grass(V,r)$, then $\Xorbit\cong K\times^{B_K}\Gorbit$. In particular, $\dim\Xorbit=\dim\Gorbit+\dim K/B_K$.
\item The correspondence preserves the closure relations. Namely, if $\Gorbit_1,\Gorbit_2$ are $B_K$-orbits of $\Grass(V,r)$ and if $\Xorbit_1=\Xi(\Gorbit_1)$ and $\Xorbit_2=\Xi(\Gorbit_2)$ are the corresponding $K$-orbits of $\Xfv$, then
\[\overline{\Gorbit_1}\subset\overline{\Gorbit_2}\iff \overline{\Xorbit_1}\subset\overline{\Xorbit_2}.\]
\end{enumerate}
\end{lemma}

Lemma \ref{L3.1} is a consequence of the following lemma (which applies to a general connected reductive group $K$).
As already used in Lemma \ref{L3.1}, $K\times^QX$ stands for the quotient of $K\times X$ by the action of $Q$ given by $q\cdot(k,x)=(kq^{-1},qx)$, and $[k,x]$ denotes the class of $(k,x)$ in this quotient. % (see \cite[5.5.8]{Springer}).

\begin{lemma}
\label{L3.2-new}
Let $K$ be a connected reductive group and let $Q\subset K$ be a parabolic subgroup. Let $X$ be an algebraic variety endowed with an action of $K$. Consider the diagonal action of $K$ on $\mathbb{X}\coloneqq K/Q\times X$.
Note that there is an isomorphism  $\chi:\mathbb{X}\to K\times^QX$ given by $\chi(kQ,x)=[k,k^{-1}x]$.
Also we consider the closed immersion $\iota:X\to\mathbb{X}$, $x\mapsto(Q,x)$.
\begin{enumerate}
\item There is a one-to-one correspondence (in fact an isomorphism of partially ordered sets)
\[\Xi:\{\mbox{$Q$-stable subsets $M\subset X$}\}\to\{\mbox{$K$-stable subsets $N\subset\mathbb{X}$}\}\]
given by $\Xi(M)=K\cdot\iota(M)$. The inverse bijection is given by $N\mapsto\iota^{-1}(N)$. Moreover, $\Xi$ restricts to a one-to-one correspondence between the orbit sets $X/Q$ and $\mathbb{X}/K$.
\item Every $Q$-stable subset $M\subset X$ yields a subset $K\times^QM\subset K\times^QX$, and we have $\chi(\Xi(M))=K\times^QM$.
\item Let $N=\Xi(M)$ for some $Q$-stable subset $M$. Then, $M$ is closed in $X$ if and only if $N$ is closed in $\mathbb{X}$. More generally, we have $\overline{N}=\Xi(\overline{M})$.
\end{enumerate}
\end{lemma}

Though this lemma is well known, we give a proof for the sake of completeness.

\begin{proof}
(1) The map $\Xi$ is certainly well defined.
For a $Q$-stable subset $M\subset X$,
the inclusion $M\subset\iota^{-1}(K\cdot\iota(M))=\iota^{-1}(\Xi(M))$ is clear, while for $x\in\iota^{-1}(\Xi(M))$ there are $k\in K$ and $y\in M$ such that $\iota(x)=k\cdot\iota(y)$, which means that $(Q,x)=(kQ,ky)$, whence $k\in Q$ and $x=ky\in Q\cdot M=M$. We have shown that $M=\iota^{-1}(\Xi(M))$.

For a $K$-stable subset $N\subset \mathbb{X}$,
given $x\in \iota^{-1}(N)$ and $q\in Q$ we have
\[\iota(qx)=(Q,qx)=q\cdot (Q,x)\in N,\]
hence $qx\in\iota^{-1}(N)$; this shows that $\iota^{-1}(N)$ is $Q$-stable.

Since $N$ is $K$-stable, the inclusion $\Xi(\iota^{-1}(N))=K\cdot \iota(\iota^{-1}(N))\subset N$ is clear, while for $(kQ,x)\in N$ we have
$\iota(k^{-1}x)=(Q,k^{-1}x)=k^{-1}\cdot(kQ,x)\in N$, hence $(kQ,x)\in K\cdot\iota(\iota^{-1}(N))=\Xi(\iota^{-1}(N))$. This shows that $N=\Xi(\iota^{-1}(N))$.

We have thus shown that $\Xi$ is a bijection, with inverse bijection given by $N\mapsto \iota^{-1}(N)$. Note also that the implications
\[(M\subset M'\Rightarrow \Xi(M)\subset\Xi(M'))\quad\mbox{and}\quad (N\subset N'\Rightarrow \iota^{-1}(N)\subset\iota^{-1}(N'))\]
are clear, which show that $\Xi$ is in fact an isomorphism of posets.

Finally, if $M=Q\cdot x$ is a $Q$-orbit, then $\Xi(M)=K\cdot(Q\cdot\iota(x))=K\cdot \iota(x)$ is a $K$-orbit. If $N=K\cdot(kQ,x)$ is a $K$-orbit, then $N=K\cdot(Q,k^{-1}x)=\Xi(Q\cdot(k^{-1}x))$ is the image of a $Q$-orbit. The proof of part (1) is complete.



(2) First we note that, if $M\subset X$ is $Q$-stable, then the quotient $K\times^QM=(K\times M)/Q$ coincides with the subset $\{[k,x]\in K\times^QX:x\in M\}\subset K\times^QX$. This subset can also be written as
\[K\times^QM=\{\chi(kQ,kx):k\in K,\ x\in M\}=\chi(K\cdot\iota(M))=\chi(\Xi(M)).\]

(3) Let $N=\Xi(M)$. If $N$ is closed, then $M=\iota^{-1}(N)$ is closed.
Conversely, assume that $M$ is closed.
Then, $\iota(M)\subset \mathbb{X}$ is closed and $Q$-stable, and this implies that
the set $\{(k,\xi)\in K\times\mathbb{X}:k^{-1}\cdot \xi\in\iota(M)\}$ is closed as well as its image in $K/Q\times\mathbb{X}$. Note that
\[N=K\cdot \iota(M)=\mathrm{pr}_2(\{(kQ,\xi)\in K/Q\times\mathbb{X}:k^{-1}\cdot \xi\in\iota(M)\}).\]
Since $K/Q$ is complete, we conclude that $N$ is closed.




More generally, using that $\Xi(\overline{M})$ and $\Xi^{-1}(\overline{N})$ are closed, and the fact that $\Xi$ is an isomorphism of posets, we get
\[\overline{N}=\overline{\Xi(M)}\subset\Xi(\overline{M})=\Xi(\overline{\Xi^{-1}(N)})\subset\Xi(\Xi^{-1}(\overline{N}))=\overline{N},\]
hence $\overline{N}=\Xi(\overline{M})$, as claimed.\end{proof}

\begin{rema}
In Lemma \ref{L3.2-new}, the assumption that $Q$ is parabolic is used only in part (3).
Also note that the isomorphism $\chi$ is $K$-equivariant when $K$ acts on $\mathbb{X}=K/Q\times X$ diagonally and on $K\times^Q X$ by left multiplication.
\end{rema}

\subsection{Parametrization and dimension formula -- proof of  Theorem \ref{Tp1}\ref{Tp11}--\ref{Tp13}}\label{section-3.2}


As before, we denote by $B_p^+\subset\GL_p(\mC)$ (resp. $B_r^+\subset\GL_r(\mC)$) the Borel subgroup of upper triangular matrices. We need two lemmas. The first one is an analogue of \cite[Proposition 6.3]{FN2}.
It is also shown in \cite[p.~390]{Fulton-1991}, but we give a proof for the sake of completeness.

\begin{lemma}
\label{L3.2}
Every $p\times r$ matrix can be written in the form
$b_1\tau b_2$
for some $b_1\in B_p^+$, $b_2\in B_r^+$, and a unique $\tau\in\ppermutationsof_{p,r}$.
\end{lemma}

\begin{proof}
Through Gauss elimination, any $p\times r$ matrix $a$ can be
transformed into a matrix $\tau\in\ppermutationsof_{p,r}$ by a
series of operations consisting of multiplying a row (resp. a
column) by a nonzero scalar or adding to a row (resp. to a column)
another row (resp. column) situated below it (resp. on its left).
These operations correspond to multiplying on the left (resp. on
the right) by an element of $B_p^+$ (resp. $B_r^+$). Hence the
double coset $B_p^+aB_r^+$ contains an element
$\tau\in\ppermutationsof_{p,r}$, which means that $a\in
B_p^+ \tau B_r^+$.

For every pair $(i,j)\in\{1,\ldots,p\}\times\{1,\ldots,r\}$, the mapping $a\mapsto \beta_{i,j}(a)\coloneqq \rank\,(a_{k,\ell})_{\substack{i\leq k\leq p\\ 1\leq \ell\leq j}}$
is constant on the set $B_p^+\tau B_r^+$, and we have
\[\beta_{i,j}(a)=\beta_{i,j}(\tau)=\card\{\mbox{$1$'s within the submatrix $(\tau_{k,\ell})_{\substack{i\leq k\leq p\\ 1\leq \ell\leq j}}$}\}.\]
This implies that two different elements $\tau,\tau'\in\ppermutationsof_{p,r}$ cannot belong to the same double coset $B_p^+aB_r^+$, whence the uniqueness.
\end{proof}

The second lemma is analogous to \cite[Lemma 8.2]{FN2}.

\begin{lemma}
\label{L3.3}
For every $\tau\in\ppermutationsof_{p,r}$, there is a permutation $w\in\permutationsof{r}$
such that $\tau w B_r^+\subset B_p^+\tau w$.
\end{lemma}

\begin{proof}
By $(e_1^\ell,\ldots,e_\ell^\ell)$ we denote the standard basis of $\mC^\ell$. By $e_{i,j}^\ell$ we denote the elementary $\ell\times \ell$ matrix
whose $(i,j)$ coefficient is $1$ and the other coefficients are $0$. By $\mathrm{diag}(t_1,\ldots,t_\ell)$, we denote the diagonal matrix with coefficients $t_1,\ldots,t_\ell$ along the diagonal.

Let $k=\rank\,\tau$. Let $1\leq i_1<\ldots<i_k\leq p$ be such that $\mathrm{Im}\,\tau=\Span{e_{i_1}^p,\ldots,e_{i_k}^p}$.
We choose $w\in\permutationsof{r}$ such that $\tau w(e_j^r)=0$ for $1\leq j\leq r-k$ and $\tau w(e_{r-k+j}^r)=e_{i_j}^p$
for $1\leq j\leq k$.

For every diagonal matrix $t=\diag(t_1,\ldots,t_r)\in B_r^+$, we have
\[\tau w t=t'\tau w\in B_p^+\tau w\]
for any diagonal matrix $t'=\diag(t'_1,\ldots,t'_p)\in B_p^+$ such that $t'_{i_j}=t_{r-k+j}$ for all $j\in\{1,\ldots,k\}$.
For every transvection $u=1_r+xe_{j,\ell}^r\in B_r^+$ where $1\leq j<\ell\leq r$, we have
\[
\tau w u=\left\{\begin{array}{ll}
\tau w & \mbox{if $j\leq r-k$},\\
(1_p+xe_{i_{j-(r-k)},i_{\ell-(r-k)}}^p)\tau w & \mbox{if $j>r-k$},
\end{array}\right.
\]
hence $\tau wu\in B_p^+\tau w$ in each case. Since $B_r^+$ is generated by the diagonal matrices and the transvections, we conclude that $\tau w B_r^+\subset B_p^+\tau w$.
\end{proof}

Now we are ready to prove parts \ref{Tp11} and \ref{Tp13} of Theorem \ref{Tp1}.


\begin{proof}[Proof of Theorem \ref{Tp1}\ref{Tp11} and \ref{Tp13}]
Every $W\in\Grass(V,r)$ is the image of a matrix
\[a=\begin{pmatrix} a_1\\ a_2 \end{pmatrix}\quad\mbox{with}\quad a_1\in\Mat_{p,r}(\mC),\ a_2\in\Mat_{q,r}(\mC),\ \rank\,a=r.\]
By Lemma \ref{L3.2}, there are $\tau_1\in \ppermutationsof_{p,r}$,
$b_1\in B_p^+$, $b_2\in B_r^+$ such that $a_1=b_1\tau_1b_2$. Moreover, by Lemma \ref{L3.3}, there is $w\in\permutationsof{r}$ such that $\tau_1wB_r^+\subset B_p^+\tau_1w$. We have
\[
a=\begin{pmatrix} a_1\\ a_2 \end{pmatrix}= \begin{pmatrix} b_1\tau_1w\\ a'_2\end{pmatrix}w^{-1}b_2
\]
for some $a'_2\in\Mat_{q,r}(\mC)$. Applying again Lemma \ref{L3.2}, there are $\tau_2\in\ppermutationsof_{q,r}$, $b_3\in B_r^+$, and $b_4\in B_q^+$ such that $a'_2=b_4\tau_2b_3$. Moreover, there is $b'_1\in B_p^+$ such that $\tau_1wb_3^{-1}=b'_1\tau_1w$. This yields
\[a=\begin{pmatrix} b_1b'_1\tau_1w\\ b_4\tau_2\end{pmatrix}b_3w^{-1}b_2=\begin{pmatrix} b_1b'_1 & 0\\ 0 & b_4\end{pmatrix}\begin{pmatrix}\tau_1w\\ \tau_2\end{pmatrix}b_3w^{-1}b_2\]
hence
\[
W=\mathrm{Im}\,a\in B_K\cdot[\omega]\quad\mbox{where}\quad \omega=\begin{pmatrix} \tau_1w\\ \tau_2\end{pmatrix}\in\ppermutationsof_{(p,q),r}.
\]
This implies that $\Grass(V,r)=\bigcup_{\omega\in\parameters}B_K\cdot[\omega]$, hence $\Xfv=\bigcup_{\omega\in\parameters}\Xorbit_\omega$ in view of Lemma \ref{L3.1}.

The mappings
\[\xi=(W,(F^+_i)_{i=0}^p,(F^-_j)_{j=0}^q)\mapsto d_{i,j}(\xi)\coloneqq \dim W\cap (F_i^++F_j^-),\] for $(i,j)\in\{0,\ldots,p\}\times\{0,\ldots,q\}$, are constant on every $K$-orbit of $\Xfv$.
Let $(e_1^+,\ldots,e_p^+)$ (resp. $(e_1^-,\ldots,e_q^-$)) be the standard basis of $V^+=\mC^p\times\{0\}^q$ (resp. $V^-=\{0\}^p\times\mC^q$), so that the standard flags $\flag_0^\pm$ are given by $\flag_0^+=(\Span{e_1^+,\ldots,e_i^+})_{i=0}^p$ and $\flag_0^-=(\Span{e_1^-,\ldots,e_j^-})_{j=0}^q$.
For every $\omega\in\parameters$,
the definition of the graph $\mathcal{G}(\omega)$ implies that the subspace \[[\omega]\cap(\Span{e_1^+,\ldots,e_i^+}+\Span{e_1^-,\ldots,e_j^-})\]
is spanned by the vectors $e_k^+$ ($1\leq k\leq i$) such that $\mathcal{G}(\omega)$ has a mark at $k^+$, the vectors $e_\ell^-$ ($1\leq \ell\leq j$) such that $\mathcal{G}(\omega)$ has a mark at $\ell^-$, and the linear combinations $e_k^++e_\ell^-$ ($1\leq k\leq i$ and $1\leq \ell\leq j$) such that $\mathcal{G}(\omega)$ has an edge joining the vertices $k^+$ and $\ell^-$. This implies that
\[
d_{i,j}(([\omega],\flag_0^+,\flag_0^-))=\dim [\omega]\cap(\Span{e_1^+,\ldots,e_i^+}+\Span{e_1^-,\ldots,e_j^-})=r_{i,j}(\omega).
\]
We deduce that
\begin{equation}
\label{3.1}
\Xorbit_\omega\subset \{\xi=(W,\flag^+,\flag^-)\in\Xfv: d_{i,j}(\xi)=r_{i,j}(\omega)\ \ \mbox{for all $i,j$}\}\quad\mbox{for all $\omega\in\parameters$}.
\end{equation}

If $\omega$, $\omega'$ are two different elements of the set $\parameters$, then their graphs $\mathcal{G}(\omega)$, $\mathcal{G}(\omega')$ must be different, hence the matrices $R(\omega)=(r_{i,j}(\omega))$ and $R(\omega')=(r_{i,j}(\omega'))$ are different. From \eqref{3.1}, it follows that the orbits $\Xorbit_\omega$ and $\Xorbit_{\omega'}$ are disjoint.
Therefore, $\Xfv$ is the disjoint union of the orbits $\Xorbit_\omega$ for $\omega\in\parameters$. This also implies that the inclusion in \eqref{3.1} must be an equality for all $\omega\in\parameters$.
This establishes Theorem \ref{Tp1}\ref{Tp11} and \ref{Tp13}.
\end{proof}

\begin{proof}[Proof of Theorem \ref{Tp1}\ref{Tp12}]
Lemma \ref{L3.1} implies
\begin{equation}
\label{dim-1}
\dim\Xorbit_\omega=\dim B_K\cdot[\omega]+\dim K/B_K=\dim B_K\cdot[\omega]+\binom{p}{2}+\binom{q}{2}.
\end{equation}
Let $\fb_K=\Lie(B_K)=\left\{\begin{pmatrix} x & 0\\ 0 & y \end{pmatrix}:x\in\fb_p^+,\ y\in\fb_q^+\right\}$, where $\fb_p^+=\Lie(B_p^+)\subset\Mat_p(\mC)$ and $\fb_q^+=\Lie(B_q^+)\subset\Mat_q(\mC)$ are the subspaces of upper triangular matrices.
We have
\begin{eqnarray}
\dim B_K\cdot[\omega] & = & \dim B_K-\dim\{b\in B_K:b([\omega])=[\omega]\}  \label{dim-2}\\
 & = & \dim\fb_K-\dim\{z\in\fb_K:z([\omega])\subset[\omega]\}. \nonumber
\end{eqnarray}
As before, we denote by $(e_1^+,\ldots,e_p^+)$, resp. $(e_1^-,\ldots,e_q^-)$, the standard basis of $V^+=\mC^p\times\{0\}^q$, resp. $V^-=\{0\}^p\times\mC^q$.
The linear space $[\omega]$ is spanned by the vectors $e_a^+$ with $a\in\{1,\ldots,p\}$ such that the graph $\mathcal{G}(\omega)$ has a mark at $a^+$, the vectors $e_c^-$ with $c\in\{1,\ldots,q\}$ such that there is a mark at $c^-$, and the linear combinations $e_a^++e_c^-$ for all pairs $(a,c)\in\{1,\ldots,p\}\times\{1,\ldots,q\}$ such that $\mathcal{G}(\omega)$ has an edge joining $a^+$ and $c^-$.
This implies that a matrix $z=\begin{pmatrix} x & 0\\ 0 & y \end{pmatrix}\in\fb_K$
satisfies $z([\omega])\subset [\omega]$ if and only if the upper triangular matrices $x=(x_{i,j})_{1\leq i,j\leq p}$ and $y=(y_{i,j})_{1\leq i,j\leq q}$ satisfy the following equations:
\begin{enumerate}
\item For every $a\in\{1,\ldots,p\}$ such that $\mathcal{G}(\omega)$ has a mark at $a^+$, we must have $x_{i,a}=0$ for all $i<a$ such that there is no mark at $i^+$.
\item For every $c\in\{1,\ldots,q\}$ such that $\mathcal{G}(\omega)$ has a mark at $c^-$, we must have $y_{j,c}=0$ for all $j<c$ such that there is no mark at $j^-$.
\item For every pair $(a,c)\in\{1,\ldots,p\}\times\{1,\ldots,q\}$ such that $\mathcal{G}(\omega)$ has an edge $(a^+,c^-)$, we must have $x_{i,a}=0$ for all $i<a$ such that $i^+$ is a free vertex (i.e. not marked nor incident with an edge) in $\mathcal{G}(\omega)$,
and we must have $y_{j,c}=0$ for all $j<c$ such that $j^-$ is a free vertex in $\mathcal{G}(\omega)$.
\item For $(a,c)$ as in (3), we must also have $x_{i,a}=y_{j,c}$ for all pair $(i,j)\in\{1,\ldots,a\}\times\{1,\ldots,c\}$ such that $(i^+,j^-)$ is an edge in $\mathcal{G}(\omega)$, i.e. for all edge which is situated on the left of $(a^+,c^-)$ or coincides with $(a^+,c^-)$ itself. Finally, we get one more equation $x_{i,a}=0$, or resp. $y_{j,c}=0$, for every edge $(i^+,j^-)$ which has a crossing with $(a^+,c^-)$, i.e. such that $i<a$ and $c<j$, resp. $i>a$ and $j<c$.
\end{enumerate}
We have listed linearly independent equations which characterize the subspace $\{z\in\fb_K:z([\omega])\subset[\omega]\}\subset\fb_K$. With the notation of Theorem \ref{Tp1}\ref{Tp12}, the above items (1)--(3) yield $a^+(\omega)+a^-(\omega)$ equations, while the item (4) yields $\frac{b(\omega)(b(\omega)+1)}{2}+c(\omega)$ equations.
This implies that
\[
\dim\{z\in\fb_K:z([\omega])\subset[\omega]\}=\dim\fb_K-\Big(a^+(\omega)+a^-(\omega)+\frac{b(\omega)(b(\omega)+1)}{2}+c(\omega)\Big).
\]
Combining this equality with \eqref{dim-1} and \eqref{dim-2}, we get the dimension formula stated in Theorem \ref{Tp1}\ref{Tp12}.
\end{proof}

\subsection{Closure relations -- proof of Theorems \ref{Tp1}\ref{Tp14} and \ref{C1}}

\label{section-3.3}

For $\omega,\omega'\in\parameters$, we write $\omega\preceq\omega'$ if we have $r_{i,j}(\omega)\geq r_{i,j}(\omega')$ for all
$(i,j)\in\{0,\ldots,p\}\times\{0,\ldots,q\}$.
This clearly endows $\parameters$ with a partial order and, for showing Theorem \ref{Tp1}\ref{Tp14}, we have to show that this order characterizes the inclusion relations between orbit closures in $\Xfv$.
We need four lemmas.

\begin{lemma}
\label{lemma-3.2}\label{L3.4}
Assume that $\omega$ is obtained from $\omega'$ by one of the elementary moves described in Figure \ref{figure1}. Then, the following relations hold:
\[\omega\prec \omega'\qquad\mbox{and}\qquad \Xorbit_\omega\subset\overline{\Xorbit_{\omega'}}.\]
\end{lemma}

\begin{proof}
We consider the cases described in Figure \ref{figure1}.
\begin{itemize}
\item
In Case 1, we have $r_{i,j}(\omega)=r_{i,j}(\omega')+1$ if $a\leq i<b$ and $c\leq j<d$, and we have $r_{i,j}(\omega)=r_{i,j}(\omega')$ otherwise.
Hence $r_{i,j}(\omega)\geq r_{i,j}(\omega')$ for all $i,j$, and this implies that $\omega\prec\omega'$.
\item In Case 2, upper subcase (resp. lower subcase), we have $r_{i,j}(\omega)=r_{i,j}(\omega')+1$ if $a\leq i<b$ and $j<c$ (resp. $i<a$ and $c\leq j<d$), and $r_{i,j}(\omega)=r_{i,j}(\omega')$ otherwise. Hence, again, we get $\omega\prec\omega'$.
\item In Case 3, upper subcase (resp. lower subcase), we have $r_{i,j}(\omega)=r_{i,j}(\omega')+1$ if $a\leq i<b$ and $c\leq j$ (resp. $a\leq i$ and $c\leq j<d$) and $r_{i,j}(\omega)=r_{i,j}(\omega')$ otherwise. Whence $\omega\prec\omega'$.
\item In Case 4, upper subcase (resp. lower subcase), we have $r_{i,j}(\omega)=r_{i,j}(\omega')+1$ if $a\leq i$ and $j<c$ (resp. $i<a$ and $c\leq j$) and $r_{i,j}(\omega)=r_{i,j}(\omega')$ otherwise. Once again this yields $\omega\prec\omega'$.
\item In Case 5, upper subcase (resp. lower subcase), we have $r_{i,j}(\omega)=r_{i,j}(\omega')+1$ if $a\leq i<b$ (resp. $c\leq j<d$) and $r_{i,j}(\omega)=r_{i,j}(\omega')$ otherwise, and once again we deduce that $\omega\prec\omega'$ in this case.
\end{itemize}
In each case, we have shown that $\omega\prec\omega'$.

As before, we denote by $e_{i,j}^k$ the $k\times k$ elementary matrix
with $1$ at position $(i,j)$ and $0$ elsewhere.
Then, let $u_{i,j}^k(t)=1_k+te_{i,j}^k$ and $\delta_i^k(t)=1_k+(t-1)e_{i,i}^k$.
For $t\in\mC^*$, we consider the matrix $h_t$ given by
\[h_t=\begin{pmatrix} A_t & 0\\ 0 & D_t \end{pmatrix},\]
where $A_t$ and $D_t$ are blocks of respective sizes $p\times p$ and $q\times q$ given by
\[
A_t=\left\{
\begin{array}{ll}
u_{a,b}^p(-t)\delta^p_a(t) & \mbox{in Cases 1, $2^+$,}\\
u_{a,b}^p(t) & \mbox{in Cases $3^+$, $5^+$,}\\
\delta_a^p(t) & \mbox{in Cases $3^-$, $4^+$,}\\
1_p & \mbox{in Cases $2^-$,\,$4^-$,\,$5^-$,}
\end{array}
\right.
\
D_t=\left\{
\begin{array}{ll}
u^q_{c,d}(t)\delta^q_c(-t) & \mbox{in Cases 1, $2^-$,}\\
u_{c,d}^q(t) & \mbox{in Cases $3^-$, $5^-$,}\\
\delta_c^q(t) & \mbox{in Cases $3^+$, $4^-$,}\\
1_q & \mbox{in Cases $2^+$,\,$4^+$,\,$5^+$.}
\end{array}
\right.
\]
Here the notation $N^+$ (resp. $N^-$) refers to the upper (resp. lower) subcase of Case $N$ in Figure \ref{figure1}.
In each case, we obtain a subset $\{h_t\}_{t\in\mC^*}\subset K$ such that
\[([\omega],\flag_0^+,\flag_0^-)=\lim_{t\to \infty}h_t\cdot([\omega'],\flag_0^+,\flag_0^-),\]
and this shows that the inclusion $\Xorbit_\omega\subset\overline{\Xorbit_{\omega'}}$ holds.
\end{proof}

\begin{lemma}
\label{L3.5}
For every $\omega,\omega'\in\parameters$, the following implication holds:
\[\Xorbit_\omega\subset\overline{\Xorbit_{\omega'}}\quad\Longrightarrow\quad \omega\preceq\omega'.\]
\end{lemma}

\begin{proof}
Assume that $\Xorbit_\omega\subset\overline{\Xorbit_{\omega'}}$. For each pair $(i,j)\in\{0,\ldots,p\}\times\{0,\ldots,q\}$, the mapping
\[
\Xfv\to\mathbb{Z}_{\geq 0},\quad (W,(F^+_k)_{k=0}^p,(F^-_\ell)_{\ell=0}^q)\mapsto \dim W\cap(F_i^++F_j^-)
\]
is upper semicontinuous. Thus, in view of Theorem \ref{Tp1}\ref{Tp13}, we have
\[
\Xorbit_\omega\subset \overline{\Xorbit_{\omega'}}\subset\{(W,(F^+_k)_{k=0}^p,(F^-_\ell)_{\ell=0}^q)\in\Xfv:\dim W\cap (F_i^++F_j^-)\geq r_{i,j}(\omega')\},
\]
whereas
\[
\Xorbit_\omega\subset\{(W,(F^+_k)_{k=0}^p,(F^-_\ell)_{\ell=0}^q)\in\Xfv:\dim W\cap (F_i^++F_j^-)=r_{i,j}(\omega)\}.
\]
This yields $r_{i,j}(\omega)\geq r_{i,j}(\omega')$ for all pair $(i,j)$, hence $\omega\preceq\omega'$.
\end{proof}

\begin{lemma}
\label{L3.6}
For every $\omega,\omega''\in\parameters$ such that $\omega\prec\omega''$, there is $\omega'\in\parameters$ with $\omega\preceq\omega'\preceq\omega''$ such that one of the pairs $(\omega,\omega')$, $(\omega',\omega'')$ fits in one of the cases described in Figure~\ref{figure1}.
\end{lemma}

\begin{proof}
We reason by induction on $p+q\geq 0$, with immediate initialization if $p+q=0$.
Assume that $\omega,\omega''\in\parameters$ are such that $\omega\prec\omega''$. In particular we have $\omega\not=\omega''$,
which forces $r\geq 1$,
i.e.  the graphs $\mathcal{G}(\omega)$ and $\mathcal{G}(\omega'')$ have at least one edge or marked vertex.

In the case where $\mathcal{G}(\omega)$ and $\mathcal{G}(\omega'')$ have
one common edge or mark -- call it $x$,
by removing this edge or mark together with the corresponding vertices (or vertex), and after renumbering of the vertices, we obtain subgraphs $\mathcal{G}(\check\omega)=\mathcal{G}(\omega)\setminus x$ and $\mathcal{G}(\check\omega'')=\mathcal{G}(\omega'')\setminus x$ associated to smaller sized matrices $\check\omega$ and $\check\omega''$, and we still have $\check\omega\prec\check\omega''$ due to the definition of the relation $\preceq$. The induction hypothesis yields $\check\omega'$ with $\check\omega\preceq\check\omega'\preceq\check\omega''$, whose associated graph $\mathcal{G}(\check\omega')$ yields $\mathcal{G}(\check\omega)$ or is yielded by $\mathcal{G}(\check\omega'')$ through one of the elementary moves described in Figure \ref{figure1}. There is an element $\omega'\in\parameters$ such that $\mathcal{G}(\check\omega')=\mathcal{G}(\omega')\setminus x$, and this element satisfies the requirements of the lemma. In conclusion,
\begin{equation}
\label{no-common-edge}
\mbox{we may assume $\mathcal{G}(\omega)$ and $\mathcal{G}(\omega'')$ have no common edge nor marked vertex.}
\end{equation}

\noindent
{\it Notation:}
It is convenient to encode the set of edges and marks of the graph $\mathcal{G}(\omega)$ in the following way:
\begin{eqnarray*}
E(\omega) & \coloneqq  & \{(a,c)\in\{1,\ldots,p\}\times\{1,\ldots,q\}:\mbox{$(a^+,c^-)$ is an edge in $\mathcal{G}(\omega)$}\}\\
 & & \cup\{(a,0):a\in\{1,\ldots,p\},\ \mbox{$\mathcal{G}(\omega)$ has a mark at $a^+$}\}\\
 & & \cup\{(0,c):c\in\{1,\ldots,q\},\ \mbox{$\mathcal{G}(\omega)$ has a mark at $c^-$}\}.
\end{eqnarray*}
Then we note that
\begin{equation}
\label{E-and-rij}
r_{i,j}(\omega)=\card E(\omega)\cap(\{0,\ldots,i\}\times\{0,\ldots,j\})\quad\mbox{for all $i,j$}.
\end{equation}
We define the set $E(\omega'')$ relative to $\omega''$ in the same way.
Both sets $E(\omega)$ and $E(\omega'')$ have $r$ elements, in particular they are nonempty.


We choose an element $(a_0,c_0)\in E(\omega'') $ with the minimal possible value of $c_0$.
If $c_0\not=0$, then $c_0^-$ is not a free vertex in $\mathcal{G}(\omega'')$: it is marked if $a_0=0$ or incident with an edge $(a_0^+,c_0^-)$ if $a_0\not=0$. Moreover, the minimality of $c_0$ guarantees then that every vertex $c^-$ with $c<c_0$ is free in $\mathcal{G}(\omega'')$, and $\mathcal{G}(\omega'')$ contains no mark at $a^+$ for all $a\in\{1,\ldots,p\}$ (because $(a,0)$ cannot belong to $E(\omega'')$, due to the minimality of $c_0$).

If $c_0=0$, then $a_0^+$ is a marked vertex in $\mathcal{G}(\omega'')$.
There may be more than one marked vertex of this type in $\mathcal{G}(\omega'')$, and we choose $a_0$ minimal for this property. Thus, in any situation, we have $r_{a_0,c_0}(\omega'')=1$.

\begin{enumerate}[label=\textbf{Case \arabic*}]
\item ($c_0\not=0$ and $r_{a_0,c_0-1}(\omega)\geq 1$)

\noindent The condition means that $\mathcal{G}(\omega)$ has an edge or mark within the vertices $\{i^+:1\leq i\leq a_0\}\cup\{j^-:1\leq j<c_0\}$.
In other words, we can find a pair $(a_1,c_1)\in E(\omega)$
with $0\leq a_1\leq a_0$ and $0\leq c_1<c_0$.
Note that $(a_0,c_1)\not=(0,0)$ since $(a_1,c_1)\not=(0,0)$.
There is an element
$\omega'\in\parameters$ such that
\[
E(\omega')=(E(\omega'')\setminus\{(a_0,c_0)\})\cup\{(a_0,c_1)\}.
\]
This incorporates several situations, and in each one the graph $\mathcal{G}(\omega')$ is deduced from $\mathcal{G}(\omega'')$ through one of the elementary moves depicted in Figure \ref{figure1}:
\begin{itemize}
\item If $a_0\not=0$ and $c_1\not=0$ (resp. $c_1=0$), then $\mathcal{G}(\omega')$ is obtained from $\mathcal{G}(\omega'')$ by replacing
the edge $(a_0^+,c_0^-)$ by an edge $(a_0^+,c_1^-)$ (resp. by a mark
at $a_0^+$), whereas $c_0^-$ becomes a free vertex. This
corresponds to Case 3 -- lower subcase (resp. Case 4 -- upper
subcase) in Figure \ref{figure1}.
\item If $a_0=0$, then $\mathcal{G}(\omega'')$
has a mark at $c_0^-$, and $\mathcal{G}(\omega')$
is obtained by replacing this mark by a mark at $c_1^-$, whereas $c_0^-$ becomes a free vertex.
This corresponds to Case 5 -- lower subcase in Figure \ref{figure1}.
\end{itemize}
In each situation, we get $\omega'\prec\omega''$ in view of Lemma \ref{lemma-3.2}.
For $(i,j)\in\{0,\ldots,p\}\times\{0,\ldots,q\}$, we have
$r_{i,j}(\omega')=r_{i,j}(\omega'')$ (hence $r_{i,j}(\omega')\leq r_{i,j}(\omega)$) unless $i\geq a_0$ and $c_1\leq j<c_0$.
If $i\geq a_0$ and $c_1\leq j<c_0$, we have
\[r_{i,j}(\omega')=r_{i,j}(\omega'')+1=1\leq r_{i,j}(\omega)\]
(the second equality is due to the minimality of $c_0$, while the inequality follows from \eqref{E-and-rij} and the fact that $(a_1,c_1)\in E(\omega)$). We conclude that the inequality $r_{i,j}(\omega')\leq r_{i,j}(\omega)$ holds for all pair $(i,j)$, hence $\omega\preceq\omega'$, and the element $\omega'$ satisfies all the requirements of the lemma.



\item ($c_0=0$ or $r_{a_0,c_0-1}(\omega)=0$)

\noindent This condition implies that the set $E(\omega)$ contains no pair of the form $(a,c)$ with $0\leq a\leq a_0$ and $0\leq c<c_0$ (see \eqref{E-and-rij}).
Note also that $(a_0,c_0)\notin E(\omega)$ (due to \eqref{no-common-edge}).

Since $r_{a_0,c_0}(\omega)\geq r_{a_0,c_0}(\omega'')=1$, there is a pair $(a'_0,c_0)\in E(\omega)$ with $0\leq a'_0< a_0$.
In particular this forces $a_0\not=0$.

The fact that $a_0\not=0$ implies that $a_0^+$ is a vertex in $\mathcal{G}(\omega)$. Either $a_0^+$ is incident with an edge/marked in $\mathcal{G}(\omega)$, in which case $E(\omega)$ contains an element of the form $(a_0,d_0)$ with $c_0< d_0\leq q$, or $a_0^+$ is a free vertex in $\mathcal{G}(\omega)$, in which case we set $d_0=q+1$.

We choose $a_1\in\{a'_0,\ldots,a_0-1\}$ maximal such that $(a_1,c_1)\in E(\omega)$ for some $c_1$ with $c_0\leq c_1<d_0$.
There is an element
$\omega'\in\parameters$ such that
\[
E(\omega')=\left\{\begin{array}{ll}
(E(\omega)\setminus\{(a_1,c_1),(a_0,d_0)\})\cup\{(a_1,d_0),(a_0,c_1)\} & \mbox{if $d_0\leq q$},\\
(E(\omega)\setminus\{(a_1,c_1)\})\cup\{(a_0,c_1)\} & \mbox{if $d_0=q+1$}.
\end{array}
\right.
\]
In each situation, the graph $\mathcal{G}(\omega')$ yields $\mathcal{G}(\omega)$ by one of the moves of Figure \ref{figure1}:
\begin{itemize}
\item In the case where $d_0\leq q$,
the graph $\mathcal{G}(\omega)$ has an edge $(a_0^+,d_0^-)$. If
$a_1,c_1\not=0$, then $(a_1^+,c_1^-)$ is also an edge in
$\mathcal{G}(\omega)$, and the relation between
$\mathcal{G}(\omega')$ and $\mathcal{G}(\omega)$ is as depicted in
Case 1 of Figure \ref{figure1}. If $c_1=0$ (resp. $a_1=0$), then
$\mathcal{G}(\omega)$ has a mark at $a_1^+$ (resp. $c_1^-$) and the
relation with $\mathcal{G}(\omega')$ is as in Case 2 - upper subcase
(resp. lower subcase) of Figure \ref{figure1}.
\item In the case where $d_0=q+1$, the vertex $a_0^+$ is a free vertex in $\mathcal{G}(\omega)$. The relation between $\mathcal{G}(\omega')$ and $\mathcal{G}(\omega)$ is as described in Case 3 - upper subcase, Case 5 - upper subcase, or Case 4 - lower subcase of Figure \ref{figure1}, depending on whether
$a_1,c_1\not=0$, $c_1=0$ (and $a_1\not=0$), or $a_1=0$ (and $c_1\not=0$).
\end{itemize}
In particular we have $\omega\prec\omega'$ (by Lemma \ref{lemma-3.2}).

For all $(i,j)\in\{0,\ldots,p\}\times\{0,\ldots,q\}$, we have $r_{i,j}(\omega')=r_{i,j}(\omega)$ unless $a_1\leq i<a_0$ and $c_1\leq j<d_0$, in which case we have $r_{i,j}(\omega')=r_{i,j}(\omega)-1$. In the latter situation, we nevertheless have
\[
r_{i,j}(\omega)=r_{a_0,j}(\omega)\geq r_{a_0,j}(\omega'')=1+r_{a_0-1,j}(\omega'')\geq 1+r_{i,j}(\omega'')
\]
(where the first equality is due to the maximality of $a_1$). Thus, the inequality $r_{i,j}(\omega')\geq r_{i,j}(\omega'')$ holds for all pair $(i,j)$, and therefore we have $\omega'\preceq\omega''$. The element $\omega'$ satisfies the required conditions. This completes the proof of the lemma.\qedhere
\end{enumerate}
\end{proof}

\begin{lemma}
\label{L3.7}
If $\Xorbit_{\omega'}$ covers $\Xorbit_\omega$, then $\dim\Xorbit_{\omega'}=\dim\Xorbit_\omega+1$.
\end{lemma}

\begin{proof}
Note that $\Xorbit_{\omega'}$ covers $\Xorbit_\omega$ if and only if $\overline{\Xorbit_\omega}$ is an irreducible component of $\overline{\Xorbit_{\omega'}}\setminus\Xorbit_{\omega'}$.
The conclusion of the lemma is implied by the following general fact, taking also Lemma \ref{L3.1} into account.

\smallskip
\noindent
{\it Fact:} Given a connected solvable algebraic group acting on an algebraic variety, the boundary $\partial O=\overline{O}\setminus O$ of each (non closed) orbit is equidimensional of codimension $1$ in $\overline{O}$.

\smallskip
\noindent
A proof of this fact can be found in \cite[Lemmas 2.12--2.13]{Timashev}.
\end{proof}

Now we are in position to proceed with the proof of Theorems \ref{Tp1}\ref{Tp14} and \ref{C1}.

\begin{proof}[Proof of Theorem \ref{Tp1}\ref{Tp14}] The ``only if'' part is shown in Lemma \ref{L3.5}. For the inverse implication, let $\omega,\omega'\in\parameters$ be such that $\omega\preceq\omega'$. Repeated applications of Lemma \ref{L3.6}
yield a sequence of elements
\[\omega=\omega_0\prec \omega_1\prec\cdots\prec \omega_\ell=\omega'\]
such that $(\omega_{k-1},\omega_k)$ fits in one of the cases of Figure \ref{figure1} for all $k$.
Then, Lemma \ref{L3.4} shows that the following sequence of inclusions holds:
\[\overline{\Xorbit_{\omega_0}}\subset \overline{\Xorbit_{\omega_1}}\subset\cdots\subset \overline{\Xorbit_{\omega_\ell}}.\]
In particular, we get the desired inclusion $\overline{\Xorbit_\omega}\subset\overline{\Xorbit_{\omega'}}$.
\end{proof}

\begin{proof}[Proof of Theorem \ref{C1}]
Assume that condition (1) of Theorem \ref{C1} holds.
First, the equality
$\dim\Xorbit_{\omega'}=\dim\Xorbit_\omega+1$
follows from Lemma \ref{L3.7}.
Next, we have in particular $\omega\prec\omega'$ in view of Theorem \ref{Tp1}\ref{Tp14}. Lemma \ref{L3.6} yields an element $\omega_0\in\parameters$ with $\omega\preceq\omega_0\preceq\omega'$ and such that $(\omega,\omega_0)$ or $(\omega_0,\omega')$ fits in one of the cases of Figure \ref{figure1}. By Theorem \ref{Tp1}\ref{Tp14} again, we get $\overline{\Xorbit_\omega}\subset\overline{\Xorbit_{\omega_0}}\subset\overline{\Xorbit_{\omega'}}$, and therefore $\omega=\omega_0$ or $\omega_0=\omega'$, due to the assumption that $\Xorbit_{\omega'}$ covers $\Xorbit_\omega$. In both cases this implies that the pair $(\omega,\omega')$ fits in one of the cases of Figure \ref{figure1}, that is, the graph $\mathcal{G}(\omega)$ is obtained from $\mathcal{G}(\omega')$ by one of the moves listed in Figure \ref{figure1}. This yields condition (2) of Theorem \ref{C1}.

Conversely, assume (2). By Lemma \ref{L3.4}, the inclusion $\overline{\Xorbit_\omega}\subset\overline{\Xorbit_{\omega'}}$ holds. This inclusion, combined with the fact that $\dim\Xorbit_{\omega'}=\dim\Xorbit_\omega+1$, implies that $\Xorbit_{\omega'}$ covers $\Xorbit_\omega$.
\end{proof}

